\documentclass[12pt,a4paper]{article}
\usepackage[UTF8]{ctex}
\usepackage{amsmath, amssymb}
\usepackage{geometry}
\usepackage{booktabs}
\usepackage{listings}
\usepackage{xcolor}
\usepackage{tcolorbox}
\usepackage{enumitem}
\usepackage{hyperref}
\usepackage{float}
\usepackage{tikz}
\usetikzlibrary{arrows.meta, positioning, shapes.geometric}

\geometry{left=2.5cm, right=2.5cm, top=2.5cm, bottom=2.5cm}

\lstset{
    language=C++,
    basicstyle=\ttfamily\small,
    keywordstyle=\color{blue}\bfseries,
    commentstyle=\color{gray},
    stringstyle=\color{red},
    showstringspaces=false,
    breaklines=true,
    frame=single,
    numbers=left,
    numberstyle=\tiny\color{gray},
    tabsize=4,
}

\title{\textbf{Q2 代码架构与运行逻辑完全指南}\\[6pt]
\large VLSI 布图规划设计 --- 固定正方形轮廓下的 HPWL 最小化}
\author{2026 年第七届"华数杯"大学生数学建模竞赛 B 题}
\date{2026/08}

\begin{document}
\maketitle
\tableofcontents
\newpage

% ============================================================
\section{项目总体架构}
% ============================================================

\subsection{目录结构}

\begin{lstlisting}[language=bash]
The_7th_Huashu_Cup.../
|-- data/raw/                    # 原始附件副本 n100/n200/n300（只读）
|-- cpp_solver/                  # C++ 共享核心（Q1/Q2/Q3 共用，单套测试）
|   `-- src/
|       |-- floor_plan.hpp       # B*-Tree 拓扑 + Skyline 解码 + 解析器
|       |-- sa.hpp               # 快速模拟退火引擎（两阶段 + 收敛 log + 快照）
|       |-- main.cpp             # CLI（mode/alpha/dead_ratio/--log/--seed）
|       |-- test_main.cpp        # 白盒测试（1135 断言）
|       `-- test_oracle.hpp      # 独立参考解码器 / 合法性校验 / 暴力枚举
|-- common/
|   |-- visualize.py             # 共享绘图基础（.rpt/.log/snap 解析）
|   `-- verify.py                # 独立合法性复核（无重叠/尺寸/越界）
|-- q2/                          # 问题 2 客制层（<-- 本文件）
|   |-- q2_solver.py             # 多起点并行求解：调 bin/main，取 HPWL 最优
|   |-- output/
|   |   |-- q2_metrics.csv       # 指标汇总（side/HPWL/利用率/合法性/seed）
|   |   |-- q2_n100.rpt          # 最优布局坐标
|   |   `-- q2_n100.log          # 收敛轨迹 + 9 等分过程快照
|   `-- visualization/
|       |-- plot_floorplan.py    # 最终布图（轮廓框 + HPWL 标注）
|       |-- plot_convergence.py  # 收敛曲线（run/run2 双阶段）
|       |-- plot_snapshots.py    # 9 等分过程快照
|       `-- figs/<时间戳>/<chip>/  # 每芯片 11 张图（保留全部历史）
|-- q1/ q3/ q4/                  # 其余问题（共享核心 + 各自客制层）
`-- paper/                       # 论文写作
\end{lstlisting}

\subsection{执行顺序}

\begin{tcolorbox}[colback=blue!5, colframe=blue!60, title=执行顺序]
\begin{enumerate}
    \item \texttt{cd cpp\_solver \&\& make}（编译 \texttt{bin/main}，0 警告）
    \item \texttt{make test}（可选，白盒测试回归，1135 断言）
    \item \texttt{conda activate math}
    \item \texttt{python q2/q2\_solver.py}（多起点并行求解 + 汇总 + 出图）
\end{enumerate}
\end{tcolorbox}

\subsection{核心设计理念}

\textbf{多起点独立退火（Multi-start SA）}：Q2 的 HPWL 目标在紧轮廓（死区 15\%）下搜索
地形崎岖，单次退火结果依赖种子。采取 N 个独立 SA 实例（不同随机种子）并行搜索、
取 HPWL 最优——这是优化领域标准做法，且天然可并行（进程级，M5 10 核吃满）。

\textbf{两阶段退火}：\texttt{run()}（自适应 $\alpha$ + 轮廓压力找可行解）$\rightarrow$
\texttt{run2()}（低温精修 HPWL）。固定轮廓下先确保合法性、再优化线长。

\textbf{可复现性}：所有轮次以 \texttt{--seed} 显式注入随机种子（\texttt{seed\_base}
记录于 CSV），结果可精确复现。

% ============================================================
\section{Q2 求解流程图}
% ============================================================

\begin{figure}[H]
\centering
\begin{tikzpicture}[
  node distance=1.1cm,
  box/.style={draw, rounded corners=2pt, fill=blue!8, align=center,
              text width=7.2cm, minimum height=0.9cm},
  dec/.style={draw, rounded corners=2pt, fill=orange!15, align=center,
              text width=7.2cm, minimum height=0.9cm},
  term/.style={draw, rounded corners=8pt, fill=green!12, align=center,
               text width=7.2cm, minimum height=0.9cm},
  arrow/.style={-{Stealth[length=3mm]}, thick},
]
\node[term] (start) {3 芯片 $\times$ N 轮并行\\（每轮独立 \texttt{--seed}，M5 10 核）};
\node[box, below=of start] (parse) {解析 .blocks/.nets/.pl\\轮廓 $W=H=\lceil\sqrt{\Sigma A \times 1.15}\rceil$};
\node[box, below=of parse] (run) {\texttt{run()} 阶段：找可行解\\自适应 $\alpha$（浮点窗口）+ 轮廓压力 $(r-R)^2$\\reset 上限兜底（单轮有界）};
\node[box, below=of run] (run2) {\texttt{run2()} 阶段：精修 HPWL\\低温冷却 + best 更新};
\node[dec, below=of run2] (filter) {可行过滤：bbox $\le$ 轮廓 side？};
\node[box, below left=1.2cm and -0.5cm of filter] (discard) {丢弃该轮};
\node[box, below right=1.2cm and -0.5cm of filter] (keep) {计入候选};
\node[box, below=of keep] (best) {N 轮中取 HPWL 最小 \texttt{best\_rpt/log}};
\node[box, below=of best] (verify) {Python 独立复核\\无重叠 / 尺寸保持 / 不越界};
\node[term, below=of verify] (out) {产出：\texttt{q2\_metrics.csv}\\ + 每芯片 11 张图（时间戳目录）};

\draw[arrow] (start) -- (parse);
\draw[arrow] (parse) -- (run);
\draw[arrow] (run) -- (run2);
\draw[arrow] (run2) -- (filter);
\draw[arrow] (filter) -- node[right] {否} (discard);
\draw[arrow] (filter) -- node[right] {是} (keep);
\draw[arrow] (discard) |- (keep);
\draw[arrow] (keep) -- (best);
\draw[arrow] (best) -- (verify);
\draw[arrow] (verify) -- (out);
\end{tikzpicture}
\caption{Q2 多起点并行求解流程图}
\end{figure}

\subsection{流程图说明}

\begin{enumerate}
    \item \textbf{多起点并行}：每个 \texttt{bin/main q2 ... --seed s} 是一个独立退火
          实例（不同初始树与扰动序列），三芯片并行、每芯片内部 N 轮并行（峰值 9 进程）。
    \item \textbf{run() 阶段}：以面积与长宽比加权代价驱动，自适应 $\alpha$ 随可行性
          窗口滑动（浮点连续化），把布局压进轮廓；reset 机制保证"找不到可行解"时
          单轮时间有界（上限 $N/16+1$ 次温度重置）。
    \item \textbf{run2() 阶段}：从 run 的 best 出发，低温精修 HPWL。
    \item \textbf{可行过滤}：bbox 超出轮廓的轮次直接丢弃（不污染最优）。
    \item \textbf{独立复核}：Python 侧重新解析文件、逐对检查重叠/尺寸/越界，
          不信任 C++ 单侧结果。
\end{enumerate}

% ============================================================
\section{Phase 1：C++ 核心 floor\_plan.hpp（Q2 视角）}
% ============================================================

\subsection{Q2 相关的关键结构}

\begin{itemize}
    \item \texttt{cost()}：返回 (HPWL, W, H)；HPWL 按引脚几何中心
          （块：$2x+w$，终端：$2x$），全量扫描线网（无增量近似）。
    \item \texttt{feas(cost)}：Q2 模式检查 \texttt{W $\le$ \_W \&\& H $\le$ \_H}。
    \item 解析器：完整读取 .blocks/.nets/.pl（HPWL 依赖终端预计算
          \texttt{NET::update}，线网按 id 排序）。
    \item \texttt{NODE} 位压缩：10-bit 父子/左右指针 + 旋转位（无符号存储，
          消除有符号位移疑虑）。
\end{itemize}

% ============================================================
\section{Phase 2：两阶段退火引擎 sa.hpp}
% ============================================================

\subsection{run()：可行性阶段}

\begin{itemize}
    \item 代价 $\text{norm\_cost} = \alpha\frac{A}{\bar A} +
          \beta\frac{\text{HPWL}}{\bar H} + (1-\alpha-\beta)\frac{(r-R)^2}{\bar r}$。
    \item $\alpha = \alpha_{base} + (1-\alpha_{base})\frac{N_{feas}}{N}$
          （浮点连续化，随可行性窗口滑动）；$\beta$ 在可行后递增施加轮廓压力。
    \item reset 机制：$tot\_feas=0$ 时温度重置，上限 $N/16+1$ 次，保证有界。
    \item 每温度层写收敛 log + 记录 best 树快照（9 等分输出）。
\end{itemize}

\subsection{run2()：精修阶段}

\begin{itemize}
    \item 从 run 的 best 出发，代价 $\text{true\_cost}$（面积与 HPWL 加权）。
    \item 初始温度 $\_init\_T/50$，低温快速下山；best 只在可行且更优时更新。
    \item \texttt{tot\_feas} 初始化为当前 best 的可行性（修复：避免从可行解
          出发仍触发 reset 空转）。
\end{itemize}

% ============================================================
\section{Phase 3：CLI main.cpp}
% ============================================================

\begin{lstlisting}[language=bash]
./bin/main q2 <alpha> <blocks> <nets> <pl> <rpt> [dead_ratio]
           [--log <file>] [--feas-only] [--seed <n>]
\end{lstlisting}

\begin{itemize}
    \item \texttt{<dead\_ratio>}：轮廓死区比例（Q2 固定 0.15；Q3 二分传不同值）。
    \item \texttt{--feas-only}：只跑 run() 判定可行性（Q3 专用，找到首可行即返回）。
    \item \texttt{--seed}：随机种子（默认当前时间），配合 \texttt{seed\_base} 复现。
    \item 输出 .rpt：cost / hpwl / area / W H / time / 每模块坐标。
\end{itemize}

% ============================================================
\section{Phase 4：多起点并行求解 q2\_solver.py}
% ============================================================

\begin{enumerate}
    \item 任务分解：\texttt{(chip, round, seed)} 全部提交线程池
          （3 芯片 $\times$ 3 轮并行 $=$ 峰值 9 进程，M5 10 核吃满）。
    \item 种子分配：\texttt{seed = seed\_base + chip\_idx*1000 + round}
          （确定性、互异、可复现）。
    \item 每轮独立输出 \texttt{q2\_<chip>\_r<round>.rpt/log}。
    \item 可行过滤（bbox $\le$ side）后取 HPWL 最小轮次为 best。
    \item 汇总 \texttt{q2\_metrics.csv}（含 \texttt{seed\_base} 复现依据）。
    \item 出图：\texttt{figs/<时间戳>/<chip>/} 每芯片 11 张
          （9 快照 + 收敛 + 布图轮廓框）。
\end{enumerate}

% ============================================================
\section{测试套件 test\_main.cpp（1135 断言）}
% ============================================================

\begin{itemize}
    \item Q2 专属用例：
          \texttt{test\_q2\_run\_snapshots}（run 阶段 9 快照格式与坐标合法）、
          \texttt{test\_q2\_feas\_only}（可行/不可行实例判定正确）、
          \texttt{test\_q2\_end\_to\_end\_hpwl}（SA 结果 HPWL 与独立 oracle 全等）。
    \item 共享 oracle 体系：独立 HPWL、暴力枚举、定义校验、穷举差分。
    \item \texttt{make test}（1135）+ \texttt{make test-san}（ASan/UBSan）双轨。
\end{itemize}

% ============================================================
\section{关键设计决策}
% ============================================================

\subsection{为什么多起点并行取最优？}

Q2 轮廓紧（死区 15\%），单次退火结果依赖种子：n300 单轮 HPWL 实测波动
574k$\sim$3152k。N 个独立实例并行、取 HPWL 最优——搜索质量单调不差于单实例，
且进程级并行吃满 10 核（实测 96-99\%）。

\subsection{为什么 \texttt{--seed} 显式化？}

默认 \texttt{srand(time)} 使结果不可复现（论文评审硬伤）。显式种子 +
\texttt{seed\_base} 记录后，任意结果可精确复现。

\subsection{为什么 run() 需要 reset 上限？}

差种子时 run() 可能长时间找不到可行解（温度重置循环）。上限
$N/16+1$ 次重置保证单轮有界（避免"卡死"），配合多起点对冲概率。

\subsection{已知弱点（如实记录，当前为冻结基线）}

\begin{itemize}
    \item run() 可行性搜索对 n300 不稳定（好/差种子单轮 1.5$\sim$7min 波动）。
    \item 8 轮取优的轮间方差大：n100 spread 192\%、n300 spread 449\%；
          中位数远差于最优（30$\sim$46\%）——best 存在运气成分。
    \item 提升方向（不改变核心算法）：增加轮数、run2 精修强度调优、
          分层筛选种子——留待参数优化阶段。
\end{itemize}

% ============================================================
\section{实际运行结果}
% ============================================================

\subsection{指标汇总（多起点 8 轮并行取优，$\alpha=0.5$，dead\_ratio=0.15）}

\begin{table}[H]
\centering
\begin{tabular}{lcccccc}
\toprule
芯片 & 轮廓 side & HPWL & 面积利用率 & 理论上限 & 合法性 \\
\midrule
n100 & 455 & 235183 & 0.867 & 0.8696 & 合法 \\
n200 & 450 & 403295 & 0.868 & 0.8696 & 合法 \\
n300 & 561 & 574322 & 0.868 & 0.8696 & 合法 \\
\bottomrule
\end{tabular}
\end{table}

面积利用率 $=$ 模块总面积 $/$ side$^2$，理论上限 $=1/1.15=0.8696$（死区 15\%）——
实测逼近理论极限，布局紧凑度最优。

\subsection{多轮统计（8 轮 HPWL 分布，评估稳定性）}

\begin{table}[H]
\centering
\begin{tabular}{lccccc}
\toprule
芯片 & min & 中位数 & max & spread & 中位/最优 \\
\midrule
n100 & 235183 & 304632 & 687833 & 192\% & 1.30 \\
n200 & 403295 & 560428 & 585763 & 45\% & 1.39 \\
n300 & 574322 & 838620 & 3152712 & 449\% & 1.46 \\
\bottomrule
\end{tabular}
\end{table}

% ============================================================
\section{文件依赖关系}
% ============================================================

\begin{lstlisting}
q2_solver.py
  |-- cpp_solver/bin/main          (编译自 main.cpp + sa.hpp + floor_plan.hpp)
  |-- data/raw/{n100,n200,n300}.{blocks,nets,pl}
  |-- common/{visualize,verify}.py
  `-- q2/visualization/plot_{floorplan,convergence,snapshots}.py
\end{lstlisting}

% ============================================================
\section{运行指南}
% ============================================================

\subsection{环境要求}

\begin{itemize}
    \item g++（Apple clang，\texttt{-std=c++17}，无第三方依赖）。
    \item conda 环境 \texttt{math}（Python 3.9+，matplotlib）。
    \item TeX Live 2026（xelatex + ctex，编译本指南）。
\end{itemize}

\subsection{完整运行流程}

\begin{lstlisting}[language=bash]
cd cpp_solver && make && make test      # 编译 + 白盒测试
cd ..
conda activate math
python q2/q2_solver.py --repeats 8      # 多起点并行求解 + 汇总 + 出图
python q2/q2_solver.py --skip-solve     # 只重出图/汇总（复用已有结果）
\end{lstlisting}

\subsection{常见问题}

\begin{enumerate}
    \item \textbf{图没更新}：求解只写 output/，画图每次新建时间戳目录——看最新目录。
    \item \textbf{HPWL 想更好}：增大 \texttt{--repeats}（纯参数，不改变算法）。
    \item \textbf{复现}：CSV 中 \texttt{seed\_base} + 每轮 \texttt{--seed} 公式
          （\texttt{seed\_base + chip\_idx*1000 + round}）即可复现。
    \item \textbf{CPU 利用}：峰值 9 进程（3 芯片 $\times$ 3 轮），M5 10 核吃满。
\end{enumerate}

% ============================================================
\section{合规清单}
% ============================================================

\begin{itemize}
    \item 原始数据 \texttt{data/raw/} 为只读副本，求解器不修改输入文件。
    \item 全部结果可通过 \texttt{make test}（1135 断言）+ \texttt{make test-san} 复现。
    \item 随机种子显式管理（\texttt{--seed} + \texttt{seed\_base}），结果可复现；
          论文数值以 \texttt{q2\_metrics.csv} 与 \texttt{figs/<时间戳>/} 为准。
    \item 已知弱点（轮间方差）如实记录，未作任何美化。
\end{itemize}

\end{document}
